\documentclass[11pt]{article}
\usepackage[margin=1.1in]{geometry}
\usepackage{amsmath, amssymb}
\usepackage{booktabs}
\usepackage{graphicx}
\usepackage[hidelinks]{hyperref}
\usepackage{caption}
\captionsetup{font=small, labelfont=bf}

\title{Auditing the Mechanism:\\
\large Ex-Ante Detection of Algorithmic Collusion from Model Internals}
\author{Briac Sockalingum\thanks{briac@berkeley.edu. First draft --- to be
completed. Method paper: ``Do Structural Models Recover Mechanisms?''
(2026), \url{https://github.com/briacSck/structural-interp}. Code and
results: \url{https://github.com/briacSck/collusion-circuits}. Hypotheses
pre-registered in \texttt{design/research\_design.md} before execution.}}
\date{July 2026 --- Working note (first draft)}

\begin{document}
\maketitle

\begin{abstract}
\noindent Competition authorities screen for cartels from market outcomes
--- prices, margins, parallelism --- after the harm has occurred. Pricing
decisions, however, are increasingly made by algorithms, and algorithms
have insides. This note asks whether a causal audit of a pricing agent's
\emph{internals}, computed from its weights and in-distribution activations
alone --- no market outcomes, no simulation of play --- can detect the
punishment mechanism that sustains algorithmic collusion. In a repeated
Bertrand laboratory calibrated to \citet{calvano2020} (our benchmarks match
theirs to three decimals), we audit neural pricing policies against the
regulator's null mechanism class --- memoryless static best response ---
by measuring the causal excess of their price responses to rival moves,
plus its undercut asymmetry. Three pre-registered findings. (i) On designed
anchor agents whose mechanism is known by construction, the audit orders
punishment structure correctly: grim trigger $>$ tit-for-tat $>$ myopic
best response $\approx$ random. (ii) As an \emph{intensity meter} the audit
fails, as we report per the pre-registration: across twenty independently
trained Q-learning pairs (all supra-competitive), the audit
\emph{anti}-correlates with the collusion index (Spearman $-0.36$) --- the
most collusive pairs sustain cooperation with credible threats that are
rarely exercised, so their expressed on-path response profiles are the
calmest. (iii) As a \emph{detector} --- the question a screening authority
actually asks --- it is strong against agents that were not trained to
evade it: the audit separates collusive from competitive mechanisms with
\textbf{AUC 0.96} against myopic controls. (iv) But a pre-registered
white-box robustness test shows the detector is \emph{trivially
Goodhart-able}: an adversarial distillation with a small audit-aware
penalty drives the audit forty-fold below the detection threshold while
leaving collusion fully intact --- the mechanism relocates to states the
audit does not measure. Mechanism screening of pricing algorithms is
feasible against today's naive agents and an arms race against strategic
ones; audit design under evasion --- randomized channels, off-path audits,
consistency cross-checks --- is where the enforcement value will be
decided. Off-path audits and the temporal-precedence question --- does the
trigger structure form \emph{before} prices rise? --- are this paper's
designated next experiments.
\end{abstract}

\section{Introduction}\label{sec:intro}

\citet{calvano2020} established that ordinary Q-learning pricing agents
learn supra-competitive prices sustained by reward--punishment schemes;
\citet{fish2024} found the same for LLM-based pricing agents; and the
empirical literature has begun documenting algorithmic pricing effects in
real markets \citep{assad2024}. The enforcement problem this creates is
well known \citep{harrington2018}: coordinated \emph{outcomes} can arise
without communication or agreement, leaving outcome-based cartel screens
--- the workhorse of proactive detection --- both legally awkward and
late. Every existing screen is behavioral and ex post.

Algorithms, unlike human cartelists, can be opened. This note develops and
stress-tests the corresponding screen: a \emph{mechanism audit}. The
regulator specifies a null mechanism class --- here, memoryless static best
response: a competitive agent's price should respond to the rival's last
price only as much as the one-shot best-response mapping does --- and the
audit measures, by causal intervention on the network's internal
activations, the \emph{excess} of the agent's price response over that
benchmark, together with its asymmetry around the cooperative price
(punishment schemes react to undercutting, not to price increases). The
audit needs the model's weights and activations at states the agent
actually visits. It needs no market data, no counterfactual simulation, and
no knowledge of whether collusion in fact occurred.

Because ``does the audit work'' is unanswerable without ground truth, the
design imports the methodology of our companion work: a population of
agents whose mechanisms are known or measured by construction. Hand-designed
\emph{anchors} (a grim-trigger network, tit-for-tat, myopic best response,
random pricing) fix the audit's validity before any learned agent is
scored; a population of independently trained Q-learning pairs, distilled
into auditable networks with 99\%+ fidelity, provides realistic learned
mechanisms spanning a range of collusive intensity; myopic-control pairs
(discount zero) provide the competitive side of the detection question.
All hypotheses were pre-registered before execution, and one of them
failed --- informatively.

\section{Laboratory and audit}\label{sec:lab}

\paragraph{Market.} Repeated Bertrand duopoly with logit demand, the
\citet{calvano2020} parameterization: our computed benchmarks
$p^N = 1.473$ and $p^M = 1.925$ match theirs; fifteen-point price grid
spanning just beyond $[p^N, p^M]$; collusion index
$\Delta = (\bar p - p^N)/(p^M - p^N)$.

\paragraph{Agents.} Tabular Q-learning pairs in self-play (memory-one
state: both firms' last prices; Calvano parameters; $10^6$ steps;
deterministic seeds), which collude reliably ($\Delta \in [0.28, 0.72]$
across twenty seeds). Each firm's greedy policy is distilled by
visitation-weighted behavior cloning into a two-hidden-layer pricing
network --- the audit target. Distillation fidelity, weighted by the states
the pair actually visits, exceeds 0.99 throughout; the distillation step is
also realistic, since deployed pricing systems are function approximators,
and it guarantees the audit target exists regardless of how the RL training
behaves. Myopic controls are identical except for a zero discount, and
converge near competitive play ($\Delta \in [0.02, 0.20]$).

\paragraph{Audit.} For the cooperative on-path row of the state space, we
patch the network's layer-level activations between adjacent rival prices,
propagate forward, and record the causal shift in the expected price. The
\emph{excess-sensitivity} channel is the positive part of this response
minus the static best-response benchmark computed from the known demand
system; the \emph{asymmetry} channel is the positive part of
(response to undercutting $-$ response to price increases). The audit
score is their sum. A caveat we state rather than bury: with memory-one
inputs, layer-level patching is close to behavioral input sensitivity;
the internal formulation earns its keep as memory lengthens, and the
deeper subspace audits of the companion paper are the designated upgrade.

\section{Results}\label{sec:results}

\subsection{Anchors (H1, confirmed): the audit orders known mechanisms}
Grim trigger 0.069 $>$ tit-for-tat 0.025 $>$ myopic best response 0.015
$\approx$ random 0.000. The ordering is the pre-registered one; the audit
reads punishment structure where it is known to exist by construction.

\subsection{Intensity (H2, rejected as pre-registered): credible threats
are quiet}
Across the twenty collusive pairs, the audit \emph{anti}-correlates with
the collusion index: Spearman $\mathbf{-0.36}$, against a pre-registered
target of $+0.6$. The mechanism, visible once seen: our v1 audit measures
the \emph{expressed} trigger structure in the on-path response profile.
The most successfully collusive pairs sustain cooperation with threats
that are credible and therefore rarely exercised --- their on-path
policies are smooth and stable --- while moderately collusive pairs cycle
through punishment episodes, leaving jagged expressed profiles. Threat
depth is an off-path object. An intensity meter must read the punishment
valleys of the value function off the equilibrium path, not the calm of
the path itself; that off-path audit is a designated next experiment of
this paper (Section~\ref{sec:next}).

\subsection{Detection (the regulator's question): AUC 0.96}
Screening authorities do not need a collusion thermometer; they need to
know which algorithms to investigate. On that question the audit performs
strongly: ranking a randomly drawn collusive pair against a randomly drawn
myopic control by audit score gives the right answer \textbf{96\%} of the
time (AUC 0.96, twenty collusive vs.\ five controls), using weights and
in-distribution activations only. Every collusive pair scores above the
myopic anchor's level. A mechanism screen of deployed pricing systems ---
prior to, or in the absence of, any market evidence --- is feasible in
this laboratory.

\subsection{Robustness (pre-registered): Goodharting the audit}
\label{sec:goodhart}
A screen that will be used will be gamed, so we attacked our own audit
under the strongest reasonable threat model: white-box, audit-aware
evasion. We re-distilled both firms of a collusive pair (base
$\Delta = 0.56$) with the loss $\mathrm{CE} + \lambda_{adv} \cdot
\mathrm{AuditProxy}$, where the proxy is the differentiable analog of the
audit, then measured the \emph{real} patching audit, self-play collusion,
and fidelity. The pre-registered evasion hypothesis was confirmed
decisively: at $\lambda_{adv} = 0.3$ the audit falls from 0.036 to
\textbf{0.001} --- forty-fold below the detection threshold --- while
self-play collusion is unchanged ($\Delta = 0.585$) and fidelity is 0.99.
The evasion sweet spot is wide ($\lambda_{adv} \in [0.3, 3]$); only at
$\lambda_{adv} \ge 10$ does suppressing the audited surface finally
destroy the collusive mechanism itself ($\Delta$ collapses to $-0.06$).
Mechanistically, the evader flattens the price-response surface exactly
where the audit measures --- the cooperative row --- while the punishment
structure that sustains cooperation survives elsewhere: the metric was
optimized, the construct relocated.

The enforcement reading is not that mechanism audits are useless --- the
AUC 0.96 result stands against agents not trained to evade, which is
today's deployed population --- but that audit design is itself a
strategic problem. Three defense directions follow directly: randomized or
undisclosed measurement channels; \emph{off-path} audits (to sustain
collusion the evader must keep punishment structure somewhere; auditing
where it must live raises the evasion cost from a regularizer to a
capability loss, as the $\lambda_{adv} \ge 10$ collapse already hints);
and consistency cross-checks between audited internals and observed
conduct. Formalizing this arms race as a mechanism-design problem is the
subject of companion theory work.

\section{Designated next experiments (this paper)}\label{sec:next}
\begin{enumerate}
\item \textbf{Off-path audits.} Read threat depth where it lives: audit
the distilled network's value/policy structure at punishment states
(post-deviation rows), not just the cooperative row. Expected to repair
the intensity meter; pre-registration to follow the standing protocol.
\item \textbf{Temporal precedence (pre-registered H3, not yet run).} Audit
mid-training checkpoints: does the trigger structure form \emph{before}
the price series turns supra-competitive? A positive answer would make the
screen genuinely ex ante in time, not only in data requirements.
\item \textbf{Inference and population size.} Permutation $p$-values and
CIs for the AUC and rank statistics; a larger population with richer
learning algorithms (DQN, policy gradient) and longer memory, where the
internal audit should genuinely dominate behavioral probing.
\end{enumerate}

\section{Discussion}\label{sec:discussion}
For enforcement: the audit reverses the screening logic --- from inferring
mechanism from outcomes to reading the mechanism and predicting outcomes
--- and its data requirements (weights, in-distribution activations) map
onto what an authority could demand in an investigation, or a marketplace
could demand of its sellers' repricers. For the interpretability program:
this is the companion methodology carried to a strategic, multi-agent
setting, with the same honesty budget --- one pre-registered hypothesis
rejected, with a mechanistic explanation and a designated repair.
Limitations: memory-one states (the behavioral-equivalence caveat above),
distilled rather than natively trained networks, one demand system, small
control group, and no legal analysis of what mechanism evidence would
support under current doctrine \citep{harrington2018}.

\begin{thebibliography}{9}\small

\bibitem[Assad et al.(2024)]{assad2024}
Assad, S., Clark, R., Ershov, D., \& Xu, L. (2024). Algorithmic pricing
and competition: Empirical evidence from the German retail gasoline
market. \emph{Journal of Political Economy Microeconomics} [TODO verify
venue/year/volume in feynman pass].

\bibitem[Calvano et al.(2020)]{calvano2020}
Calvano, E., Calzolari, G., Denicol\`o, V., \& Pastorello, S. (2020).
Artificial intelligence, algorithmic pricing, and collusion.
\emph{American Economic Review}, 110(10), 3267--3297.

\bibitem[Fish et al.(2024)]{fish2024}
Fish, S., Gonczarowski, Y.~A., \& Shorrer, R.~I. (2024). Algorithmic
collusion by large language models. \emph{arXiv:2404.00806} [TODO verify
arXiv id in feynman pass].

\bibitem[Harrington(2018)]{harrington2018}
Harrington, J.~E. (2018). Developing competition law for collusion by
autonomous artificial agents. \emph{Journal of Competition Law \&
Economics}, 14(3), 331--363 [TODO verify pages in feynman pass].

\bibitem[Sockalingum(2026)]{companion2026}
Sockalingum, B. (2026). Do structural models recover mechanisms? Evidence
from neural agents with observable internals. Working note,
\url{https://github.com/briacSck/structural-interp}.

\end{thebibliography}

% TODO (author + feynman pass): verify Assad et al., Fish et al., and
% Harrington citation metadata (flagged inline); add Klein (2021, RAND) if
% verified; AUC inference; figure population.png; DG COMP-facing summary
% paragraph if used for the motivation letter.

\end{document}
